\cvsection{Education}

\begin{center}
  \textbf{Software Engineering}, 2015 | Gewerblich-Technische Schulen Offenburg, Offenburg im Kinzigtal, Germany\\
  \textbf{Realschule}, 2012 | Realschule Schramberg, Schramberg, Germany
\end{center}

\cvsection{Language Proficiencies}

\begin{center}
  German (Native) | English (Proficient)
\end{center}

\cvsection{Extra-Curricular Activities}

\jobheading{DevOps-Lead Kubernetes rc3.world}{Chaos Computer Club}{12/2021}

Supported the rc3.world backend team as part of the rC3 infrastructure team, overseeing and maintaining a Kubernetes cluster and a complex multi-staged Elixir deployment hosting the remote congress experience.

\begin{itemize}
  \cvitemplain{Orchestrated a highly available Elixir Application}{using Kubernetes, Kustomize, Helm, ArgoCD, ingress-nginx, cert-manager, and the kube-prometheus-stack.}
  \cvitemplain{Optimised build \& deployment process}{by enhancing the Dockerfile layout to use a multi-stage build process, ensuring a smaller and more efficient production image.}
  \cvitemplain{Built and maintained a multi-stage GitLab CI pipeline}{to compile, verify, and deploy documentation for the Elixir application powering the event.}
\end{itemize}

\jobheading{Head of Network, GPN18/19}{Entropia e.V.}{01/2018 to 06/2019}

Led the network infrastructure team of over 15 volunteers with a diverse background for one of the largest CCC events in Germany after the Chaos Communication Congress, acting as the primary representative towards event management, handling sponsor acquisition, multi-homed BGP connectivity, budget management, and logistics for a fault-tolerant network.

\cvsection{Personal Projects}

\begin{itemize}
  \cvitem{Home Lab}{Operate two Kubernetes clusters using ClusterAPI on Hetzner, K3s on Raspberry Pis, as well as a Raspberry Pi based CEPH storage cluster, providing a solid testbed for distributed systems research.}
  \cvitem{Deployment \& Automation}{ArgoCD/GitOps-based Kubernetes deployments and Ansible for bare metal services.}
  \cvitem{Observability \& Monitoring}{Maintaining a full observability platform running on Kubernetes consisting of the Grafana LGTM Stack (Mimir, Loki, and Tempo) and Thanos for metrics, logs, and traces, ensuring real-time visibility into performance.}
  \cvitem{Networking \& Security}{Build and manage a zero-trust home network using Tailscale, leveraging ACL-based access control for network segmentation and security.}
  \cvitem{Open Source Contributions}{Actively contribute to multiple open-source projects.}
\end{itemize}

\cvsection{Core Skills}

\begin{itemize}
  \item Go, Kubernetes Operators \& Distributed Systems
  \item Multi-Cloud Infrastructure (AWS, Azure, Google Cloud, OCI)
  \item Chaos Engineering \& Resilience Testing
  \item CI/CD, Deployment Orchestration, Infrastructure as Code
  \item Low-Level Systems Programming
  \item Technical Leadership, Mentorship \& Stakeholder Management
\end{itemize}

\cvsection{Talks}

\begin{itemize}
  \item \href{https://media.ccc.de/v/gpn21-48-site-reliability-engineering-explained-an-exploration-of-devops-platform-engineering-and-sre}{Site Reliability Engineering Explained - An Exploration of DevOps, Platform Engineering, and SRE}
  \item \href{https://media.ccc.de/v/gpn20-22-understanding-alerting-how-to-come-up-with-a-good-enough-alerting-strategy}{Understanding Alerting - How to come up with a good enough alerting strategy}
  \item \href{https://media.ccc.de/v/gpn21-47-modern-observability-scalable-observability-with-the-lgtm-stack-harnessing-the-power-of-loki-grafana-tempo-and-mimir}{Modern Observability - Scalable Observability with the LGTM Stack: Harnessing the Power of Loki, Grafana, Tempo, and Mimir}
  \item \href{https://media.ccc.de/v/gpn21-49-from-0-to-kubernetes-eine-einfhrung-zur-container-orchestrierung-mit-praktischen-antworten-auf-die-hufigsten-fragen-wie-warum-oder-wann-}{From 0 to Kubernetes Eine Einführung zur Container-Orchestrierung mit praktischen Antworten auf die häufigsten Fragen wie ``warum?'' oder ``wann?''}
  \item \href{https://media.ccc.de/v/gpn22-393-congratulations-it-s-a-kubernetes-but-now-what-}{Congratulations, it's a Kubernetes! But now what?}
  \item \href{https://media.ccc.de/v/denog14-22895-hilfe-ich-bin-manager-und-jetzt-oder-auch-wie-man-in-der-it-karriere-machen-kann-}{Hilfe, ich bin Manager - und jetzt? Oder auch: Wie man in der IT Karriere machen kann}
\end{itemize}
